\section{Methodology}

% The goal of this paper is to answer the following question: "To what extent should insurance companies be allowed to use statistical profiling to determine policy eligibility and pricing?" Since this question does not have a simple answer, subquestions will be answered so that a conclusion to the main question can be drawn. The following subquestions will be used:
% \begin{enumerate}
% 	\item What types of statistical data do insurance companies typically use in risk assessment?
% 	\item Do any of the statistics used raise concerns about discrimination and if so, who are they discriminant against?
% 	\item ?
% 	\item ?
% \end{enumerate}
%
% This risk assessment will be examined through two frameworks: an ethical and a political framework. Within the ethical framework, two main approaches will be used: deontology and consequentialism.
%
% Deontology is an ethical view that determines the morality of an action based on whether the action itself is moral. ``Kant believed moral agents were on duty to comply with what he believed to be a universally valid moral law. According to Kant moral rules are thus supposed to be universal laws. Like the law of gravitational attraction, it is universal and applies to all people equally regardless of status, need or identity''~\cite{syllabus}. This framework is important, because insurance companies make moral choices about fairness in risk assessment. It can be argued that the models that are used unfairly deny insurance based on factors like socioeconomic status, race, gender, etc.
%
% Consequentialism is an ethical view that states that the only thing that matters is the outcome and not the morality of the action itself~\cite{syllabus}. This is a contradictory view to deontology, as one focuses on the action itself and the other just on the outcome. This view is useful in determining whether the results of statistical profiling are harmful or not.

The question of whether companies should be allowed to vary pricing based of statistical profiling is a very broad one. Naturally, we will analyse the issue by looking at various subsets of the risk factors involved for profiling. Such subsets can be defined based (as seen in~\cite{odihi20}) based on criteria like whether the agent is responsible for the risk factor, whether the risk factor is causally linked to health costs, and whether the agent can be reasonably expected to know about the aforementioned causal link.

Before we proceed, we will outline a few examples of risk factors commonly used for statistical profiling that fall on the different sides of each criteria above. For the first criteria, age is a risk factor the agent is not responsible for, while their level of physical activity is one they usually are. For the second criteria, people who smoke are more likely to attain certain diseases (todo: add source), which lead to higher medical costs (todo: insert source), thus whether the agent smokes is a risk factor that is causally responsible for the increased health costs. On the other hand, while a person's place of living is could be used for profiling if correlation was observed with health costs, without requiring the causality link to exist (as mentioned in~\cite{odihi20}). As for the last criteria, while people in developed countries can, in the 21st century, be reasonably expected to know about the health risks of smoking (insert source), this would not have been the case a hundred years ago (insert source).

\subsection{The ethical framework}
We will analyse the ramifications of the different risk factors through two frameworks: an ethical and a policy-based one. For our ethical analysis, we will consider the ethical theories of consequentialism and deontology.

In particular, we will use utilitarianism (a form of consequentialism). According to~\cite[Chapter~7]{elemofmp}, utilitarianism is based around the ``Principle of Utility'', that is, the principle that all actions should attempt to maximize the so called ``utility'' produced. Intuitively, our analysis through the lens of utilitarianism will focus on the consequences profiling using different classes of risk factors has on the parties involved.

Second, we will analyse the issues through the lens of deontology. Kant's ethics (i.e.~deontology) are based around the idea that humans should ``always be treated as an end and never as a means only'', which he formulate into a categorical imperative (a law of morality, which all actions must abide by): ``Act so that you treat humanity, whether in your own person or in that of another, always as an end and never as a means only''~\cite[Chapter~10]{elemofmp}. A more concrete equivalent formulation of this imperative, as stated by~\cite[Section~4.2.1]{syllabus} is the following:~''Act only according to that maxim whereby you can, at the same time, will that it should become a universal law''.
% TODO: add one more sentence explaining what this means for our paper, practically

\subsection{The policy framework}
According to~\cite{syllabus}, ``policy is a preferred course of action, with the desired consequence of solving a problem''.

